\mxunit{premiere-automatismes}{Banque d’automatismes}\label{AUTO}
Deux questions par séance, en 6 à 8 minutes correction comprise. Les identifiants permettent des reprises ; ils ne constituent pas une liste à effectuer intégralement à chaque cours. Les dates de disponibilité sont fondées sur l’enseignement prévu, à décaler si la notion n’a pas encore été étudiée.
\mxneed{5}\subsection*{R001 -- Fractions}\label{R001}
Calculer $\frac23+\frac14$.
\ifmxteacher\textbf{Correction.} $\dfrac{11}{12}$ : dénominateur commun 12.\par\textit{Prérequis :} Acquis antérieurs\par\textit{Disponible à partir du 2026-09-03 ; reprises conseillées :} 2026-09-10, 2026-09-24, 2026-11-03.\fi
\mxneed{5}\subsection*{R002 -- Priorités}\label{R002}
Calculer $5-2(3-7)$.
\ifmxteacher\textbf{Correction.} $13$ : parenthèses, produit, différence.\par\textit{Prérequis :} Acquis antérieurs\par\textit{Disponible à partir du 2026-09-03 ; reprises conseillées :} 2026-09-10, 2026-09-24, 2026-11-03.\fi
\mxneed{5}\subsection*{R003 -- Puissances}\label{R003}
Simplifier $\dfrac{10^5\times 10^{-3}}{10^4}$.
\ifmxteacher\textbf{Correction.} $10^{-2}=\num{0,01}$.\par\textit{Prérequis :} Acquis antérieurs\par\textit{Disponible à partir du 2026-09-03 ; reprises conseillées :} 2026-09-10, 2026-09-24, 2026-11-03.\fi
\mxneed{5}\subsection*{R004 -- Racines}\label{R004}
Simplifier $\sqrt{50}$.
\ifmxteacher\textbf{Correction.} $5\sqrt2$.\par\textit{Prérequis :} Acquis antérieurs\par\textit{Disponible à partir du 2026-09-03 ; reprises conseillées :} 2026-09-10, 2026-09-24, 2026-11-03.\fi
\mxneed{5}\subsection*{R005 -- Proportion}\label{R005}
18 représente \pc{30} d’un total. Quel total ?
\ifmxteacher\textbf{Correction.} $\dfrac{18}{\num{0,3}}=60$.\par\textit{Prérequis :} Acquis antérieurs\par\textit{Disponible à partir du 2026-09-03 ; reprises conseillées :} 2026-09-10, 2026-09-24, 2026-11-03.\fi
\mxneed{5}\subsection*{R006 -- Taux direct}\label{R006}
Un prix de 80 baisse de \pc{15}. Valeur finale ?
\ifmxteacher\textbf{Correction.} $80\times\num{0,85}=68$.\par\textit{Prérequis :} Acquis antérieurs\par\textit{Disponible à partir du 2026-09-03 ; reprises conseillées :} 2026-09-10, 2026-09-24, 2026-11-03.\fi
\mxneed{5}\subsection*{R007 -- Valeur initiale}\label{R007}
Après une hausse de \pc{20}, la valeur est 72. Valeur initiale ?
\ifmxteacher\textbf{Correction.} $\dfrac{72}{\num{1,2}}=60$.\par\textit{Prérequis :} Acquis antérieurs\par\textit{Disponible à partir du 2026-09-03 ; reprises conseillées :} 2026-09-10, 2026-09-24, 2026-11-03.\fi
\mxneed{5}\subsection*{R008 -- Taux calculé}\label{R008}
Une quantité passe de 40 à 50. Taux ?
\ifmxteacher\textbf{Correction.} $\dfrac{50-40}{40}=\pc{25}$.\par\textit{Prérequis :} Acquis antérieurs\par\textit{Disponible à partir du 2026-09-03 ; reprises conseillées :} 2026-09-10, 2026-09-24, 2026-11-03.\fi
\mxneed{5}\subsection*{R009 -- Taux successifs}\label{R009}
Hausse de \pc{20} puis baisse de \pc{20} : taux global ?
\ifmxteacher\textbf{Correction.} $\num{1,2}\times\num{0,8}=\num{0,96}$ : baisse de \pc{4}.\par\textit{Prérequis :} Acquis antérieurs\par\textit{Disponible à partir du 2026-09-03 ; reprises conseillées :} 2026-09-10, 2026-09-24, 2026-11-03.\fi
\mxneed{5}\subsection*{R010 -- Taux réciproque}\label{R010}
Quel taux compense une baisse de \pc{20} ?
\ifmxteacher\textbf{Correction.} $\dfrac{1}{\num{0,8}}=\num{1,25}$ : hausse de \pc{25}.\par\textit{Prérequis :} Acquis antérieurs\par\textit{Disponible à partir du 2026-09-03 ; reprises conseillées :} 2026-09-10, 2026-09-24, 2026-11-03.\fi
\mxneed{5}\subsection*{R011 -- Produit nul}\label{R011}
Résoudre $(2x-1)(x+3)=0$.
\ifmxteacher\textbf{Correction.} $x=\dfrac{1}{2}$ ou $x=-3$.\par\textit{Prérequis :} Acquis antérieurs\par\textit{Disponible à partir du 2026-09-03 ; reprises conseillées :} 2026-09-10, 2026-09-24, 2026-11-03.\fi
\mxneed{5}\subsection*{R012 -- Développement}\label{R012}
Développer $(2x-3)^2$.
\ifmxteacher\textbf{Correction.} $4x^2-12x+9$.\par\textit{Prérequis :} Acquis antérieurs\par\textit{Disponible à partir du 2026-09-03 ; reprises conseillées :} 2026-09-10, 2026-09-24, 2026-11-03.\fi
\mxneed{5}\subsection*{R013 -- Factorisation}\label{R013}
Factoriser $3x^2-12x$.
\ifmxteacher\textbf{Correction.} $3x(x-4)$.\par\textit{Prérequis :} Acquis antérieurs\par\textit{Disponible à partir du 2026-09-03 ; reprises conseillées :} 2026-09-10, 2026-09-24, 2026-11-03.\fi
\mxneed{5}\subsection*{R014 -- Signe affine}\label{R014}
Résoudre $-3x+6\geq0$.
\ifmxteacher\textbf{Correction.} $x\leq2$ : inverser l’inégalité à la division par -3.\par\textit{Prérequis :} Acquis antérieurs\par\textit{Disponible à partir du 2026-09-03 ; reprises conseillées :} 2026-09-10, 2026-09-24, 2026-11-03.\fi
\mxneed{5}\subsection*{R015 -- Pente}\label{R015}
Pente par $\pt{1;4}$ et $\pt{3;0}$ ?
\ifmxteacher\textbf{Correction.} $-2$.\par\textit{Prérequis :} Acquis antérieurs\par\textit{Disponible à partir du 2026-09-03 ; reprises conseillées :} 2026-09-10, 2026-09-24, 2026-11-03.\fi
\mxneed{5}\subsection*{R016 -- Droite}\label{R016}
Équation de pente -2 par $\pt{1;4}$ ?
\ifmxteacher\textbf{Correction.} $y=-2x+6$.\par\textit{Prérequis :} Acquis antérieurs\par\textit{Disponible à partir du 2026-09-03 ; reprises conseillées :} 2026-09-10, 2026-09-24, 2026-11-03.\fi
\mxneed{5}\subsection*{R017 -- Moyenne pondérée}\label{R017}
Les valeurs 2 et 5 ont effectifs 3 et 2. Moyenne ?
\ifmxteacher\textbf{Correction.} $\dfrac{6+10}{5}=\num{3,2}$.\par\textit{Prérequis :} Acquis antérieurs\par\textit{Disponible à partir du 2026-09-03 ; reprises conseillées :} 2026-09-10, 2026-09-24, 2026-11-03.\fi
\mxneed{5}\subsection*{R018 -- Médiane}\label{R018}
Médiane de 1, 2, 2, 5, 8, 10 ?
\ifmxteacher\textbf{Correction.} $\dfrac{2+5}{2}=\num{3,5}$ selon la convention moyenne des deux valeurs centrales.\par\textit{Prérequis :} Acquis antérieurs\par\textit{Disponible à partir du 2026-09-03 ; reprises conseillées :} 2026-09-10, 2026-09-24, 2026-11-03.\fi
\mxneed{5}\subsection*{R019 -- Quartiles}\label{R019}
Pour la série 1, 2, 3, 4, 5, 6, 7, 8, donner $Q_1$ et $Q_3$ avec les rangs plafonds $\left\lceil\dfrac n4\right\rceil$ et $\left\lceil\dfrac{3n}4\right\rceil$.
\ifmxteacher\textbf{Correction.} $Q_1=2$, $Q_3=6$.\par\textit{Prérequis :} Acquis antérieurs\par\textit{Disponible à partir du 2026-09-03 ; reprises conseillées :} 2026-09-10, 2026-09-24, 2026-11-03.\fi
\mxneed{5}\subsection*{R020 -- Écart type statistique}\label{R020}
Pour 1 et 3 équipondérés, moyenne, variance et écart type ?
\ifmxteacher\textbf{Correction.} Moyenne 2, variance $\dfrac{[(1-2)^2+(3-2)^2]}{2}=1$, écart type 1.\par\textit{Prérequis :} Acquis antérieurs\par\textit{Disponible à partir du 2026-09-03 ; reprises conseillées :} 2026-09-10, 2026-09-24, 2026-11-03.\fi
\mxneed{5}\subsection*{R021 -- Unités}\label{R021}
Convertir 1 min 40 s en secondes ; 2400 images sur cette durée donnent combien par seconde ?
\ifmxteacher\textbf{Correction.} $100$ s ; $24$ images/s.\par\textit{Prérequis :} Acquis antérieurs\par\textit{Disponible à partir du 2026-09-03 ; reprises conseillées :} 2026-09-10, 2026-09-24, 2026-11-03.\fi
\mxneed{5}\subsection*{R022 -- Probabilité simple}\label{R022}
Sur un dé équilibré, probabilité d’un multiple de 3 ?
\ifmxteacher\textbf{Correction.} $\dfrac{2}{6}=\dfrac{1}{3}$.\par\textit{Prérequis :} Acquis antérieurs\par\textit{Disponible à partir du 2026-09-03 ; reprises conseillées :} 2026-09-10, 2026-09-24, 2026-11-03.\fi
\mxneed{5}\subsection*{R023 -- Complémentaire}\label{R023}
$P(A)=\dfrac{3}{8}$. Donner $P(\overline A)$.
\ifmxteacher\textbf{Correction.} $\dfrac{5}{8}$.\par\textit{Prérequis :} Acquis antérieurs\par\textit{Disponible à partir du 2026-09-03 ; reprises conseillées :} 2026-09-10, 2026-09-24, 2026-11-03.\fi
\mxneed{5}\subsection*{R024 -- Intervalle}\label{R024}
Écrire $x\geq-2$ et $x<3$ en intervalle.
\ifmxteacher\textbf{Correction.} $\intfo{-2}{3}$.\par\textit{Prérequis :} Acquis antérieurs\par\textit{Disponible à partir du 2026-09-03 ; reprises conseillées :} 2026-09-10, 2026-09-24, 2026-11-03.\fi
\mxneed{5}\subsection*{R025 -- Logique}\label{R025}
Nier « tous les élèves ont réussi ».
\ifmxteacher\textbf{Correction.} « Au moins un élève n’a pas réussi. »\par\textit{Prérequis :} Acquis antérieurs\par\textit{Disponible à partir du 2026-09-03 ; reprises conseillées :} 2026-09-10, 2026-09-24, 2026-11-03.\fi
\mxneed{5}\subsection*{R026 -- Histogramme}\label{R026}
Deux classes de largeurs 2 et 4 ont le même effectif 20. Dans un histogramme à aires proportionnelles aux effectifs, quelles hauteurs relatives ?
\ifmxteacher\textbf{Correction.} Hauteurs 10 et 5 : l’aire, pas la hauteur seule, code l’effectif.\par\textit{Prérequis :} Acquis antérieurs\par\textit{Disponible à partir du 2026-09-03 ; reprises conseillées :} 2026-09-10, 2026-09-24, 2026-11-03.\fi
\mxneed{5}\subsection*{R027 -- Diagramme circulaire}\label{R027}
Un secteur mesure $72^\circ$. Proportion ?
\ifmxteacher\textbf{Correction.} $\dfrac{72}{360}=\pc{20}$.\par\textit{Prérequis :} Acquis antérieurs\par\textit{Disponible à partir du 2026-09-03 ; reprises conseillées :} 2026-09-10, 2026-09-24, 2026-11-03.\fi
\mxneed{5}\subsection*{R028 -- Échelle}\label{R028}
Un axe commence à 90. Des barres s’arrêtent à 100 et 110. La seconde valeur est-elle le double de la première ?
\ifmxteacher\textbf{Correction.} Non : $\dfrac{110}{100}=\num{1,1}$. Les hauteurs dessinées au-dessus de 90 trompent cette comparaison.\par\textit{Prérequis :} Acquis antérieurs\par\textit{Disponible à partir du 2026-09-03 ; reprises conseillées :} 2026-09-10, 2026-09-24, 2026-11-03.\fi
\mxneed{5}\subsection*{R029 -- Pourcentage de pourcentage}\label{R029}
\pc{40} des élèves font A ; parmi eux, \pc{25} font B. Proportion qui fait A et B ?
\ifmxteacher\textbf{Correction.} $\num{0,4}\times\num{0,25}=\pc{10}$.\par\textit{Prérequis :} Acquis antérieurs\par\textit{Disponible à partir du 2026-09-03 ; reprises conseillées :} 2026-09-10, 2026-09-24, 2026-11-03.\fi
\mxneed{5}\subsection*{R030 -- Proportion et effectif}\label{R030}
Une catégorie représente \num{0,125} de 80 personnes. Effectif ?
\ifmxteacher\textbf{Correction.} $10$ personnes.\par\textit{Prérequis :} Acquis antérieurs\par\textit{Disponible à partir du 2026-09-03 ; reprises conseillées :} 2026-09-10, 2026-09-24, 2026-11-03.\fi
\mxneed{5}\subsection*{R031 -- Définir et représenter une suite}\label{R031}
Calculer les quatre premiers termes de $w_0=2$, $w_{n+1}=w_n+2n+1$ ; proposer une formule explicite.
\ifmxteacher\textbf{Correction.} $2,3,6,11$. La formule $w_n=n^2+2$ convient : elle donne le premier terme et $(n+1)^2+2-(n^2+2)=2n+1$.\par\textit{Prérequis :} Calcul littéral, puissances entières, pourcentages, fonction affine, boucle Python.\par\textit{Disponible à partir du 2026-09-09 ; reprises conseillées :} 2026-09-17, 2026-10-01, 2026-11-03.\fi
\mxneed{5}\subsection*{R032 -- Suites arithmétiques}\label{R032}
On sait $u_3=17$ et $u_8=37$ pour une suite arithmétique. Déterminer $r$ et $u_0$.
\ifmxteacher\textbf{Correction.} $37-17=5r$, donc $r=4$ puis $u_0=17-3\times4=5$.\par\textit{Prérequis :} Calcul littéral, puissances entières, pourcentages, fonction affine, boucle Python.\par\textit{Disponible à partir du 2026-09-11 ; reprises conseillées :} 2026-09-22, 2026-10-06, 2026-11-03.\fi
\mxneed{5}\subsection*{R033 -- Suites géométriques}\label{R033}
Étudier la variation de $u_n=-2\times3^n$, puis de $v_n=2(-1)^n$.
\ifmxteacher\textbf{Correction.} $u_{n+1}-u_n=-4\times3^n<0$ : décroissante. $v_n$ alterne 2 et $-2$ : ni croissante ni décroissante.\par\textit{Prérequis :} Calcul littéral, puissances entières, pourcentages, fonction affine, boucle Python.\par\textit{Disponible à partir du 2026-09-16 ; reprises conseillées :} 2026-09-24, 2026-10-08, 2026-11-05.\fi
\mxneed{5}\subsection*{R034 -- Sommes, indices et démonstrations}\label{R034}
Calculer $\sum_{k=2}^{5}2^k$ sans symbole de sommation dans la réponse.
\ifmxteacher\textbf{Correction.} $4+8+16+32=60$, ou $\dfrac{4(1-2^4)}{1-2}=60$.\par\textit{Prérequis :} Calcul littéral, puissances entières, pourcentages, fonction affine, boucle Python.\par\textit{Disponible à partir du 2026-09-18 ; reprises conseillées :} 2026-09-29, 2026-10-13, 2026-11-10.\fi
\mxneed{5}\subsection*{R035 -- Conjecturer une limite et programmer un seuil}\label{R035}
Pour $u_n=48(\dfrac{1}{2})^n$, déterminer les premiers rangs où $u_n\leq6$ et où $u_n<6$.
\ifmxteacher\textbf{Correction.} $u_3=6$ : rang 3 pour $\leq6$, rang 4 pour $<6$.\par\textit{Prérequis :} Calcul littéral, puissances entières, pourcentages, fonction affine, boucle Python.\par\textit{Disponible à partir du 2026-09-18 ; reprises conseillées :} 2026-09-29, 2026-10-13, 2026-11-10.\fi
\mxneed{5}\subsection*{R036 -- Trois formes, trois lectures}\label{R036}
Donner trois formes de $g(x)=x^2+2x-3$. Lire ses racines et son minimum.
\ifmxteacher\textbf{Correction.} $g(x)=(x+1)^2-4=(x-1)(x+3)$. Racines 1 et $-3$, minimum $-4$ en $-1$.\par\textit{Prérequis :} Identités remarquables, produit nul, signe d’un produit, racines carrées.\par\textit{Disponible à partir du 2026-09-23 ; reprises conseillées :} 2026-10-01, 2026-10-15, 2026-11-12.\fi
\mxneed{5}\subsection*{R037 -- Discriminant et résolution}\label{R037}
Résoudre $3x^2-6x+3=0$ et $x^2+2x+5=0$.
\ifmxteacher\textbf{Correction.} La première est $3(x-1)^2=0$, donc 1. Pour la seconde, $\Delta=-16<0$ : aucune solution réelle.\par\textit{Prérequis :} Identités remarquables, produit nul, signe d’un produit, racines carrées.\par\textit{Disponible à partir du 2026-09-25 ; reprises conseillées :} 2026-10-06, 2026-11-03, 2026-11-17.\fi
\mxneed{5}\subsection*{R038 -- Factoriser sans automatiser le discriminant}\label{R038}
Déterminer le polynôme f de degré 2 tel que $f(1)=f(4)=0$ et $f(0)=-8$.
\ifmxteacher\textbf{Correction.} $f(x)=a(x-1)(x-4)$ ; $4a=-8$, donc $f(x)=-2(x-1)(x-4)$.\par\textit{Prérequis :} Identités remarquables, produit nul, signe d’un produit, racines carrées.\par\textit{Disponible à partir du 2026-09-30 ; reprises conseillées :} 2026-10-08, 2026-11-03, 2026-11-19.\fi
\mxneed{5}\subsection*{R039 -- Signe, inéquations et optimisation}\label{R039}
Résoudre $x^2-4x-5<0$ puis maximiser $x(8-x)$ sur $\intff{0}{3}$.
\ifmxteacher\textbf{Correction.} $(x-5)(x+1)<0$ donne $\intoo{-1}{5}$. Le sommet de $x(8-x)$ est en 4, hors du domaine ; sur $\intff{0}{3}$, la fonction croît et son maximum vaut 15 en 3.\par\textit{Prérequis :} Identités remarquables, produit nul, signe d’un produit, racines carrées.\par\textit{Disponible à partir du 2026-10-02 ; reprises conseillées :} 2026-10-13, 2026-11-03, 2026-11-24.\fi
\mxneed{5}\subsection*{R040 -- Projection et angle}\label{R040}
Deux vecteurs ont pour normes 2 et 5 et forment $90^\circ$, puis $180^\circ$. Calculer les produits.
\ifmxteacher\textbf{Correction.} $0$ puis $-10$. Le cosinus vaut respectivement 0 et $-1$.\par\textit{Prérequis :} Vecteurs, Chasles, Pythagore, cosinus dans un triangle rectangle, repère orthonormé.\par\textit{Disponible à partir du 2026-10-07 ; reprises conseillées :} 2026-10-15, 2026-11-03, 2026-11-26.\fi
\mxneed{5}\subsection*{R041 -- Coordonnées, norme et orthogonalité}\label{R041}
Pour quel réel t les vecteurs $\pt{t;2}$ et $\pt{3;-6}$ sont-ils orthogonaux ?
\ifmxteacher\textbf{Correction.} $3t-12=0$, donc $t=4$.\par\textit{Prérequis :} Vecteurs, Chasles, Pythagore, cosinus dans un triangle rectangle, repère orthonormé.\par\textit{Disponible à partir du 2026-10-09 ; reprises conseillées :} 2026-11-03, 2026-11-03, 2026-12-01.\fi
\mxneed{5}\subsection*{R042 -- Calcul vectoriel et identités}\label{R042}
On donne $\norm{\vec u}=2$, $\norm{\vec v}=3$, $\vec u\cdot\vec v=1$. Calculer $\norm{\vec u-\vec v}$.
\ifmxteacher\textbf{Correction.} Le carré vaut $4-2+9=11$ ; la norme vaut $\sqrt{11}$.\par\textit{Prérequis :} Vecteurs, Chasles, Pythagore, cosinus dans un triangle rectangle, repère orthonormé.\par\textit{Disponible à partir du 2026-10-14 ; reprises conseillées :} 2026-11-03, 2026-11-05, 2026-12-03.\fi
\mxneed{5}\subsection*{R043 -- Al-Kashi et choix d’une méthode}\label{R043}
Un triangle a $AB=AC=4$ et $BC=4\sqrt2$. Déterminer l’angle A.
\ifmxteacher\textbf{Correction.} $32=16+16-32\cos A$, donc $\cos A=0$ et $A=90^\circ$.\par\textit{Prérequis :} Vecteurs, Chasles, Pythagore, cosinus dans un triangle rectangle, repère orthonormé.\par\textit{Disponible à partir du 2026-10-16 ; reprises conseillées :} 2026-11-03, 2026-11-10, 2026-12-08.\fi
\mxneed{5}\subsection*{R044 -- Le cercle de diamètre AB}\label{R044}
Pour $A\pt{0;0}$, $B\pt{6;0}$, déterminer les M tels que $\vect{MA}\cdot\vect{MB}=7$.
\ifmxteacher\textbf{Correction.} $MI^2-9=7$ : cercle de centre $\pt{3;0}$ et de rayon 4.\par\textit{Prérequis :} Vecteurs, Chasles, Pythagore, cosinus dans un triangle rectangle, repère orthonormé.\par\textit{Disponible à partir du 2026-10-16 ; reprises conseillées :} 2026-11-03, 2026-11-10, 2026-12-08.\fi
\mxneed{5}\subsection*{R045 -- Taux de variation et sécantes}\label{R045}
Calculer le taux de $f(x)=x^2+3x$ entre 2 et $2+h$.
\ifmxteacher\textbf{Correction.} $f(2+h)-f(2)=7h+h^2$ ; taux $7+h$ pour $h\ne0$.\par\textit{Prérequis :} Images, pente d’une droite, calcul littéral, racine carrée et valeur absolue.\par\textit{Disponible à partir du 2026-11-04 ; reprises conseillées :} 2026-11-12, 2026-11-26, 2027-01-05.\fi
\mxneed{5}\subsection*{R046 -- Nombre dérivé}\label{R046}
Pour $f(x)=3x^2-2x$, calculer $f\prime(1)$ à partir du taux.
\ifmxteacher\textbf{Correction.} Le taux vaut $4+3h$, donc le nombre dérivé vaut 4.\par\textit{Prérequis :} Images, pente d’une droite, calcul littéral, racine carrée et valeur absolue.\par\textit{Disponible à partir du 2026-11-06 ; reprises conseillées :} 2026-11-17, 2026-12-01, 2027-01-05.\fi
\mxneed{5}\subsection*{R047 -- Tangente et approximation affine}\label{R047}
On donne $f(3)=5$ et $f\prime(3)=-2$. Écrire la tangente et approcher $f(\num{3,04})$.
\ifmxteacher\textbf{Correction.} $y=5-2(x-3)=-2x+11$ ; approximation $5-2(\num{0,04})=\num{4,92}$.\par\textit{Prérequis :} Images, pente d’une droite, calcul littéral, racine carrée et valeur absolue.\par\textit{Disponible à partir du 2026-11-11 ; reprises conseillées :} 2026-11-19, 2026-12-03, 2027-01-05.\fi
\mxneed{5}\subsection*{R048 -- Quand la dérivabilité échoue}\label{R048}
Justifier que $g(x)=|x-2|$ n’est pas dérivable en 2. Étudier aussi la dérivabilité de $x\mapsto\sqrt{x}$ en 0 par son taux.
\ifmxteacher\textbf{Correction.} Le taux en 2 est $\dfrac{|h|}{h}$, avec les deux valeurs limites opposées. Pour la racine en 0, le taux $\dfrac{1}{\sqrt h}$ pour $h>0$ dépasse toute borne lorsque h se rapproche de 0 : aucun nombre dérivé réel.\par\textit{Prérequis :} Images, pente d’une droite, calcul littéral, racine carrée et valeur absolue.\par\textit{Disponible à partir du 2026-11-13 ; reprises conseillées :} 2026-11-24, 2026-12-08, 2027-01-05.\fi
\mxneed{5}\subsection*{R049 -- Conditionner : changer le dénominateur}\label{R049}
Sur 60 objets, 20 sont rouges ; parmi les rouges, 5 sont abîmés. Calculer $P(R\cap D)$ et $P_R(D)$.
\ifmxteacher\textbf{Correction.} $\dfrac{5}{60}=\dfrac{1}{12}$ et $\dfrac{5}{20}=\dfrac{1}{4}$.\par\textit{Prérequis :} Fractions, proportions, probabilités sur un univers fini, intersection et réunion.\par\textit{Disponible à partir du 2026-11-18 ; reprises conseillées :} 2026-11-26, 2026-12-10, 2027-01-07.\fi
\mxneed{5}\subsection*{R050 -- Arbres, partitions et probabilités totales}\label{R050}
$P(A)=\dfrac{1}{3}$, $P_A(B)=\dfrac{3}{4}$, $P_{\overline A}(B)=\dfrac{1}{4}$. Calculer $P(B)$.
\ifmxteacher\textbf{Correction.} $\dfrac{1}{4}+\dfrac{1}{6}=\dfrac{5}{12}$.\par\textit{Prérequis :} Fractions, proportions, probabilités sur un univers fini, intersection et réunion.\par\textit{Disponible à partir du 2026-11-20 ; reprises conseillées :} 2026-12-01, 2026-12-15, 2027-01-12.\fi
\mxneed{5}\subsection*{R051 -- Indépendance et incompatibilité}\label{R051}
On donne $P(A)=\num{0,4}$, $P(B)=\num{0,5}$, $P(A\cap B)=\num{0,18}$. Sont-ils indépendants ?
\ifmxteacher\textbf{Correction.} Non : $\num{0,4}\times\num{0,5}=\num{0,20}\ne\num{0,18}$.\par\textit{Prérequis :} Fractions, proportions, probabilités sur un univers fini, intersection et réunion.\par\textit{Disponible à partir du 2026-11-25 ; reprises conseillées :} 2026-12-03, 2026-12-17, 2027-01-19.\fi
\mxneed{5}\subsection*{R052 -- Épreuves successives et Bernoulli}\label{R052}
Quatre essais indépendants de probabilité $p=\dfrac{1}{3}$. Calculer la probabilité de quatre échecs puis d’au moins un succès.
\ifmxteacher\textbf{Correction.} $(\dfrac{2}{3})^4=\dfrac{16}{81}$ puis $1-\dfrac{16}{81}=\dfrac{65}{81}$.\par\textit{Prérequis :} Fractions, proportions, probabilités sur un univers fini, intersection et réunion.\par\textit{Disponible à partir du 2026-11-27 ; reprises conseillées :} 2026-12-08, 2027-01-05, 2027-01-19.\fi
\mxneed{5}\subsection*{R053 -- Simuler et discuter le modèle}\label{R053}
Une simulation de 2000 points donne 675 points sous $y=x^2$. Quelle aire estime-t-elle ?
\ifmxteacher\textbf{Correction.} $\dfrac{675}{2000}=\num{0,3375}$ ; c’est une estimation, pas la valeur exacte.\par\textit{Prérequis :} Fractions, proportions, probabilités sur un univers fini, intersection et réunion.\par\textit{Disponible à partir du 2026-11-27 ; reprises conseillées :} 2026-12-08, 2027-01-05, 2027-01-19.\fi
\mxneed{5}\subsection*{R054 -- Fonction dérivée et fonctions usuelles}\label{R054}
Dériver $x^5$, $x^{-2}$ et $\sqrt x$ en précisant les domaines.
\ifmxteacher\textbf{Correction.} $5x^4$ sur $\R$ ; $-\dfrac{2}{x^3}$ pour $x\ne0$ ; $\dfrac{1}{2\sqrt x}$ pour $x>0$.\par\textit{Prérequis :} Nombre dérivé, calcul algébrique, signe d’un trinôme, intervalles.\par\textit{Disponible à partir du 2026-12-02 ; reprises conseillées :} 2026-12-10, 2027-01-05, 2027-01-21.\fi
\mxneed{5}\subsection*{R055 -- Somme, produit, inverse et quotient}\label{R055}
Dériver $(x^2+1)(3x-2)$ puis $\dfrac{1}{x^2+1}$.
\ifmxteacher\textbf{Correction.} $2x(3x-2)+3(x^2+1)=9x^2-4x+3$ ; $-\dfrac{2x}{(x^2+1)^2}$, sur $\R$.\par\textit{Prérequis :} Nombre dérivé, calcul algébrique, signe d’un trinôme, intervalles.\par\textit{Disponible à partir du 2026-12-04 ; reprises conseillées :} 2026-12-15, 2027-01-05, 2027-01-26.\fi
\mxneed{5}\subsection*{R056 -- Signe de la dérivée et extrema}\label{R056}
Étudier $f(x)=x^2-4x+1$ sur $\intff{0}{3}$ et donner ses extrema sur cet intervalle.
\ifmxteacher\textbf{Correction.} $f\prime=2x-4$, donc décroissance jusqu’à 2 puis croissance. Minimum $-3$ en 2 ; valeurs aux bornes 1 et $-2$, donc maximum 1 en 0.\par\textit{Prérequis :} Nombre dérivé, calcul algébrique, signe d’un trinôme, intervalles.\par\textit{Disponible à partir du 2026-12-09 ; reprises conseillées :} 2026-12-17, 2027-01-05, 2027-01-28.\fi
\mxneed{5}\subsection*{R057 -- Parité, inégalités et comparaison}\label{R057}
Étudier la parité de $x^4+x^2$ et $x^3-x$ ; démontrer $x^2+1\geq2x$.
\ifmxteacher\textbf{Correction.} Première paire, seconde impaire ; différence $x^2+1-2x=(x-1)^2\geq0$, égalité en 1.\par\textit{Prérequis :} Nombre dérivé, calcul algébrique, signe d’un trinôme, intervalles.\par\textit{Disponible à partir du 2026-12-11 ; reprises conseillées :} 2027-01-05, 2027-01-05, 2027-02-02.\fi
\mxneed{5}\subsection*{R058 -- Optimisation}\label{R058}
Un rectangle de périmètre 12 m a largeur x. Déterminer l’aire maximale par dérivation.
\ifmxteacher\textbf{Correction.} $A(x)=x(6-x)$ sur $\intff{0}{6}$ ; $A\prime=6-2x$. Maximum 9 m$^2$ en 3.\par\textit{Prérequis :} Nombre dérivé, calcul algébrique, signe d’un trinôme, intervalles.\par\textit{Disponible à partir du 2026-12-11 ; reprises conseillées :} 2027-01-05, 2027-01-05, 2027-02-02.\fi
\mxneed{5}\subsection*{R059 -- Vecteur normal et équation cartésienne}\label{R059}
Équation de la droite par $\pt{2;-1}$ de normale $\pt{3;4}$.
\ifmxteacher\textbf{Correction.} $3(x-2)+4(y+1)=0$, soit $3x+4y-2=0$.\par\textit{Prérequis :} Coordonnées vectorielles, produit scalaire, systèmes linéaires, complétion du carré. Repère toujours orthonormé.\par\textit{Disponible à partir du 2027-01-06 ; reprises conseillées :} 2027-01-19, 2027-01-28, 2027-03-02.\fi
\mxneed{5}\subsection*{R060 -- Projeté orthogonal et distance}\label{R060}
Projeter $P\pt{4;1}$ sur $d:x=2$.
\ifmxteacher\textbf{Correction.} $H\pt{2;1}$ et distance 2 ; la normale est horizontale.\par\textit{Prérequis :} Coordonnées vectorielles, produit scalaire, systèmes linéaires, complétion du carré. Repère toujours orthonormé.\par\textit{Disponible à partir du 2027-01-08 ; reprises conseillées :} 2027-01-19, 2027-02-02, 2027-03-02.\fi
\mxneed{5}\subsection*{R061 -- Équation d’un cercle}\label{R061}
Reconnaître $x^2+y^2+2x-8y+17=0$.
\ifmxteacher\textbf{Correction.} $(x+1)^2+(y-4)^2=0$ : le seul point $\pt{-1;4}$, pas un cercle de rayon strictement positif.\par\textit{Prérequis :} Coordonnées vectorielles, produit scalaire, systèmes linéaires, complétion du carré. Repère toujours orthonormé.\par\textit{Disponible à partir du 2027-01-13 ; reprises conseillées :} 2027-01-21, 2027-02-04, 2027-03-04.\fi
\mxneed{5}\subsection*{R062 -- Résoudre une configuration}\label{R062}
Intersecter $(x-1)^2+(y+2)^2=9$ avec $y=1$.
\ifmxteacher\textbf{Correction.} $(x-1)^2=0$, donc seul point $\pt{1;1}$ ; la droite est tangente au cercle.\par\textit{Prérequis :} Coordonnées vectorielles, produit scalaire, systèmes linéaires, complétion du carré. Repère toujours orthonormé.\par\textit{Disponible à partir du 2027-01-15 ; reprises conseillées :} 2027-01-26, 2027-02-09, 2027-03-09.\fi
\mxneed{5}\subsection*{R063 -- Variable, événement et loi}\label{R063}
Deux pièces équilibrées sont lancées. X compte les piles. Écrire sa loi et $P(X\leq1)$.
\ifmxteacher\textbf{Correction.} Valeurs 0,1,2 ; probabilités $\dfrac{1}{4}$,$\dfrac{1}{2}$,$\dfrac{1}{4}$. $P(X\leq1)=\dfrac{3}{4}$.\par\textit{Prérequis :} Probabilités, moyenne pondérée, carrés, arbres et listes.\par\textit{Disponible à partir du 2027-01-20 ; reprises conseillées :} 2027-01-28, 2027-02-11, 2027-03-11.\fi
\mxneed{5}\subsection*{R064 -- Espérance et linéarité}\label{R064}
Un gain brut vaut 10 avec probabilité $\dfrac{1}{5}$ et 0 sinon. Quelle mise rend le jeu équitable ?
\ifmxteacher\textbf{Correction.} $E(G)=\dfrac{10}{5}=2$. Pour une mise m, $E(G-m)=2-m$ ; mise équitable 2.\par\textit{Prérequis :} Probabilités, moyenne pondérée, carrés, arbres et listes.\par\textit{Disponible à partir du 2027-01-22 ; reprises conseillées :} 2027-02-02, 2027-03-02, 2027-03-16.\fi
\mxneed{5}\subsection*{R065 -- Variance, écart type et König–Huygens}\label{R065}
Pour X de valeurs $-1,1$ équiprobables, calculer espérance, variance et écart type.
\ifmxteacher\textbf{Correction.} $E(X)=0$, $E(X^2)=1$, variance 1 et écart type 1.\par\textit{Prérequis :} Probabilités, moyenne pondérée, carrés, arbres et listes.\par\textit{Disponible à partir du 2027-01-27 ; reprises conseillées :} 2027-02-04, 2027-03-02, 2027-03-18.\fi
\mxneed{5}\subsection*{R066 -- Transformer et comparer des gains}\label{R066}
$E(X)=3$, $V(X)=4$. Donner les caractéristiques de $Y=-2X+5$.
\ifmxteacher\textbf{Correction.} $E(Y)=-1$, $V(Y)=16$, $\sigma(Y)=4$.\par\textit{Prérequis :} Probabilités, moyenne pondérée, carrés, arbres et listes.\par\textit{Disponible à partir du 2027-01-29 ; reprises conseillées :} 2027-02-09, 2027-03-02, 2027-03-23.\fi
\mxneed{5}\subsection*{R067 -- Échantillons et moyenne observée}\label{R067}
X prend les valeurs 0 et 2 de façon équiprobable. Calculer la moyenne observée de l’échantillon 2, 0, 2, 2, 0, puis son écart à $E(X)$.
\ifmxteacher\textbf{Correction.} $m=\dfrac{6}{5}=\num{1,2}$ et $E(X)=1$ : l’écart vaut \num{0,2} pour cet échantillon.\par\textit{Prérequis :} Probabilités, moyenne pondérée, carrés, arbres et listes.\par\textit{Disponible à partir du 2027-01-29 ; reprises conseillées :} 2027-02-09, 2027-03-02, 2027-03-23.\fi
\mxneed{5}\subsection*{R068 -- Définir une fonction nouvelle}\label{R068}
Quelle est l’équation de la tangente à $e^x$ en 1 ?
\ifmxteacher\textbf{Correction.} $y=e+e(x-1)=ex$.\par\textit{Prérequis :} Puissances, suites géométriques, dérivation et signe.\par\textit{Disponible à partir du 2027-02-03 ; reprises conseillées :} 2027-02-11, 2027-03-02, 2027-03-25.\fi
\mxneed{5}\subsection*{R069 -- Propriétés algébriques et comparaisons}\label{R069}
Simplifier $\dfrac{e^{3x}}{e^{x-1}}$ ; résoudre $e^{x+2}>e^{2x}$.
\ifmxteacher\textbf{Correction.} $e^{2x+1}$ ; $x+2>2x$, donc $x<2$.\par\textit{Prérequis :} Puissances, suites géométriques, dérivation et signe.\par\textit{Disponible à partir du 2027-02-05 ; reprises conseillées :} 2027-03-02, 2027-03-02, 2027-03-30.\fi
\mxneed{5}\subsection*{R070 -- Dériver et étudier}\label{R070}
Dériver $g(t)=3e^{-2t}$ et étudier son sens de variation.
\ifmxteacher\textbf{Correction.} $g\prime(t)=-6e^{-2t}<0$ sur $\R$ : strictement décroissante.\par\textit{Prérequis :} Puissances, suites géométriques, dérivation et signe.\par\textit{Disponible à partir du 2027-02-10 ; reprises conseillées :} 2027-03-02, 2027-03-04, 2027-04-06.\fi
\mxneed{5}\subsection*{R071 -- Du discret au continu}\label{R071}
Pour $Q(t)=5e^{\num{0,3}t}$, exprimer $Q(n)$ et $\dfrac{Q(t+2)}{Q(t)}$.
\ifmxteacher\textbf{Correction.} $u_n=5(e^{\num{0,3}})^n$ ; le rapport sur deux unités vaut $e^{\num{0,6}}$.\par\textit{Prérequis :} Puissances, suites géométriques, dérivation et signe.\par\textit{Disponible à partir du 2027-02-12 ; reprises conseillées :} 2027-03-02, 2027-03-09, 2027-04-06.\fi
\mxneed{5}\subsection*{R072 -- Construire numériquement et démontrer une inégalité}\label{R072}
Montrer que $e^{2x}\geq1+2x$ pour tout réel x.
\ifmxteacher\textbf{Correction.} Appliquer l’inégalité démontrée avec l’argument $2x$ ; égalité seulement en x=0.\par\textit{Prérequis :} Puissances, suites géométriques, dérivation et signe.\par\textit{Disponible à partir du 2027-02-12 ; reprises conseillées :} 2027-03-02, 2027-03-09, 2027-04-06.\fi
\mxneed{5}\subsection*{R073 -- Mesurer un angle en radians}\label{R073}
Convertir $45^\circ$, $120^\circ$ ; donner l’arc de rayon 2 m et angle $\dfrac{3\pi}{4}$.
\ifmxteacher\textbf{Correction.} $\dfrac{\pi}{4}$, $\dfrac{2\pi}{3}$ ; longueur $\dfrac{3\pi}{2}$ m.\par\textit{Prérequis :} Cercle, Pythagore, cosinus et sinus dans un triangle rectangle, symétries, fractions.\par\textit{Disponible à partir du 2027-03-03 ; reprises conseillées :} 2027-03-11, 2027-03-25, 2027-04-27.\fi
\mxneed{5}\subsection*{R074 -- Enroulement et image d’un réel}\label{R074}
Placer les images de $\dfrac{7\pi}{2}$ et $-\dfrac{5\pi}{3}$ en les ramenant à un tour.
\ifmxteacher\textbf{Correction.} $\dfrac{7\pi}{2}$ a l’image de $\dfrac{3\pi}{2}$ ; $-\dfrac{5\pi}{3}$ celle de $\dfrac{\pi}{3}$.\par\textit{Prérequis :} Cercle, Pythagore, cosinus et sinus dans un triangle rectangle, symétries, fractions.\par\textit{Disponible à partir du 2027-03-05 ; reprises conseillées :} 2027-03-16, 2027-03-30, 2027-04-27.\fi
\mxneed{5}\subsection*{R075 -- Sinus, cosinus et valeurs remarquables}\label{R075}
Si $\cos x=\dfrac{3}{5}$ et l’image de x est dans le quatrième quadrant, donner $\sin x$.
\ifmxteacher\textbf{Correction.} $\sin ^2x=\dfrac{16}{25}$ ; le sinus est négatif, donc $-\dfrac{4}{5}$.\par\textit{Prérequis :} Cercle, Pythagore, cosinus et sinus dans un triangle rectangle, symétries, fractions.\par\textit{Disponible à partir du 2027-03-10 ; reprises conseillées :} 2027-03-18, 2027-04-06, 2027-04-29.\fi
\mxneed{5}\subsection*{R076 -- Angles associés et symétries}\label{R076}
Calculer les deux coordonnées des images de $\dfrac{4\pi}{3}$ et $\dfrac{7\pi}{4}$.
\ifmxteacher\textbf{Correction.} $\pt{-\dfrac{1}{2};-\dfrac{\sqrt 3}{2}}$ et $\pt{\dfrac{\sqrt 2}{2};-\dfrac{\sqrt 2}{2}}$.\par\textit{Prérequis :} Cercle, Pythagore, cosinus et sinus dans un triangle rectangle, symétries, fractions.\par\textit{Disponible à partir du 2027-03-12 ; reprises conseillées :} 2027-03-23, 2027-04-06, 2027-05-04.\fi
\mxneed{5}\subsection*{R077 -- Approximer pi et réinvestir}\label{R077}
Pour n=6 et c=1, calculer la nouvelle corde après doublement.
\ifmxteacher\textbf{Correction.} $h=\dfrac{\sqrt 3}{2}$ ; nouvelle corde $\dfrac{1}{\sqrt {2+\sqrt3}}=\sqrt {2-\sqrt3}$.\par\textit{Prérequis :} Cercle, Pythagore, cosinus et sinus dans un triangle rectangle, symétries, fractions.\par\textit{Disponible à partir du 2027-03-12 ; reprises conseillées :} 2027-03-23, 2027-04-06, 2027-05-04.\fi
\mxneed{20}
\mxneed{5}\subsection*{R078 -- Lecture graphique}\label{R078}
Sur la figure R078, lire $f(1)$, résoudre $f(x)=0$ et $f(x)>0$.
\begin{center}\begin{tikzpicture}\begin{axis}[width=.8\linewidth,height=5cm,axis lines=middle,xmin=-2.5,xmax=4.5,ymin=-2,ymax=3,xtick={-2,-1,0,1,2,3,4},ytick={-1,0,1,2},grid=both]\addplot[very thick,mxmint,mark=*] coordinates {(-2,-1)(1,2)(4,-1)};\end{axis}\end{tikzpicture}\end{center}\ifmxteacher\textbf{Correction.} $f(1)=2$ ; zéros $-1$ et 3 ; signe positif sur $\intoo{-1}{3}$.\par\textit{Prérequis :} Repérage\par\textit{Disponible à partir du 2026-09-03 ; reprises conseillées :} 2026-09-10, 2026-09-24, 2026-11-03.\fi
